
\section{Agentic Consensus: The Case for a Primary Engineering Artifact}

Human control over AI-generated systems requires a primary artifact that encodes structural commitments. Code encodes only execution; the consensus layer $C$ fills this gap.

\subsection{Structural Artifact Foundations}

\subsubsection{$C$ as the Coordinating Artifact}
We model a project state as a tuple $(I, C, A, E)$:
\begin{itemize}[leftmargin=*, itemsep=2pt]
  \item $I$: user intent, potentially partial, implicit, and time-varying,
  \item $C$: a structured consensus layer (schema, graph, views),
  \item $A$: realized artifacts that execute in the world (code, configs, dataflows), and
  \item $E$: evidence linked to claims in $C$ (tests, traces, provenance).
\end{itemize}
Although the four components are defined together, they are not peers: $C$ is the management layer through which the other three become coherent.
$I$ enters the system through propose/clarify operations that encode intent into structural claims within $C$.
$A$ is derived from $C$ via the realization operator $\Phi$; without $C$, code exists but lacks an auditable rationale.
$E$ attaches to claims in $C$ as validation signals; without $C$, test results are disconnected observations rather than evidence for or against specific commitments.
In short, without $C$ the remaining components are individually meaningful but collectively uncoordinated and unauditable.

\subsubsection{$C$ as an Auditable Graph}
The consensus layer is not a single diagram; it is a typed, queryable schema with multiple synchronized projections.
In practice, $C$ is a labeled property graph whose nodes represent modules, interfaces, data entities, and invariants; edges represent dependencies, contracts, and temporal relations.

%% \begin{figure*}[t]
%%   \centering
%%   \begin{tikzpicture}[
%%     node distance=6mm and 18mm,
%%     nd-human/.style={
%%       draw=black!70, fill=black!6, rounded corners=4pt,
%%       inner sep=6pt, align=center, font=\small,
%%       minimum width=32mm, minimum height=10mm
%%     },
%%     nd-center/.style={
%%       draw=black!85, fill=black!14, rounded corners=4pt,
%%       inner sep=7pt, align=center, font=\small\bfseries,
%%       minimum width=46mm, minimum height=14mm, thick
%%     },
%%     nd-artifact/.style={
%%       draw=black!70, fill=black!6, rounded corners=4pt,
%%       inner sep=6pt, align=center, font=\small,
%%       minimum width=32mm, minimum height=10mm
%%     },
%%     nd-evidence/.style={
%%       draw=black!55, fill=black!3, rounded corners=3pt,
%%       inner sep=5pt, align=center, font=\footnotesize,
%%       minimum width=32mm, minimum height=9mm, dashed
%%     },
%%     nd-views/.style={
%%       draw=black!55, fill=black!3, rounded corners=3pt,
%%       inner sep=5pt, align=center, font=\footnotesize,
%%       minimum width=46mm, minimum height=9mm, dashed
%%     },
%%     ar-main/.style={-{Latex[length=2.5mm, width=2mm]}, thick},
%%     ar-back/.style={-{Latex[length=2.5mm, width=2mm]}, thick, dashed},
%%     ar-evid/.style={-{Latex[length=2mm, width=1.5mm]}, thin, draw=black!50},
%%     nd-lbl/.style={font=\footnotesize, inner sep=2pt, fill=white}
%%   ]
%%
%%   %% ── Row 1: Intent | Consensus Layer | Artifacts ──
%%   \node[nd-human]    (intent)
%%         {Intent $I$\\\footnotesize{(natural language, pointing)}};
%%   \node[nd-center, right=of intent] (C)
%%         {Consensus Layer $C$\\\normalfont\footnotesize{typed property graph,}\\
%%          \footnotesize{schema + synchronized views}};
%%   \node[nd-artifact, right=of C] (artifacts)
%%         {Artifacts $A$\\\footnotesize{code, configs, dataflows}};
%%
%%   %% ── Row 2: Views under C | Evidence under A ──
%%   \node[nd-views, below=of C] (views)
%%         {Multi-modal projections\\\footnotesize{topology, lineage, experiments,}\\
%%          \footnotesize{contracts, causal paths}};
%%   \node[nd-evidence, below=of artifacts] (evidence)
%%         {Evidence $E$\\\footnotesize{tests, checks, traces, provenance}};
%%
%%   %% ── Main horizontal arrows ──
%%   \draw[ar-main] (intent) -- node[nd-lbl, above]{propose / clarify} (C);
%%   \draw[ar-main] (C)      -- node[nd-lbl, above]{$\Phi$\;{\normalfont\footnotesize realize}} (artifacts);
%%   \draw[ar-back] (artifacts) to [bend left=28]
%%                  node[nd-lbl, below]{$\Psi$\;{\normalfont\footnotesize rehydrate}} (C);
%%
%%   %% ── Vertical arrows ──
%%   \draw[ar-evid] (C)         -- (views);
%%   \draw[ar-evid] (artifacts) -- (evidence);
%%   \draw[ar-evid] (evidence.west) to [bend right=20]
%%                  node[nd-lbl, left]{\footnotesize link} (C.south east);
%%
%%   %% ── Human review feedback loop ──
%%   \draw[ar-back] (C.north west) to [bend left=35]
%%                  node[nd-lbl, above]{\footnotesize review / commit / rollback} (intent.north east);
%%
%%   %% ── Dashed boundary around C + views (the "Consensus Zone") ──
%%   \node[draw=black!55, dashed, rounded corners=6pt,
%%         fit=(C)(views), inner sep=8pt,
%%         label={[font=\footnotesize, text=black!60]above right:Consensus Zone}] {};
%%
%%   %% ── Distinguishing annotation ──
%%   \node[font=\footnotesize, text=black!60, align=left,
%%         below=5mm of views, xshift=16mm]
%%         {\textit{Differs from UML / notebooks:}\quad
%%          (i)~round-trip $\Phi$/$\Psi$ operators\quad
%%          (ii)~evidence-linked claims\quad
%%          (iii)~entropy-driven clarification};
%%
%%   \end{tikzpicture}
%%   \caption{%
%%     Architecture of Agentic Consensus.
%%     The consensus layer $C$ (bold center) is the primary artifact;
%%     executable artifacts $A$ are derived via $\Phi$ (realize) and
%%     kept in round-trip correspondence via $\Psi$ (rehydrate).
%%     Evidence $E$ links directly to structural claims in $C$, making every
%%     commitment auditable.
%%     Humans review and commit at the level of $C$, not raw diffs.%
%%   }
%%   \label{fig:canvas}
%% \end{figure*}

\begin{figure*}[t]
  \centering
  % \includegraphics[width=\textwidth]{figures/figure-cases.pdf}
  \includegraphics[width=0.90\textwidth]{figures/figure-cases.pdf}
  \caption{Two case studies contrasting vibe coding (left sub-panels) with Agentic Consensus (right sub-panels). \textbf{Case~1} (left): an ML training pipeline's silent AUC degradation---vibe coding produces scattered, context-free debugging sessions, while $C$ maintains pipeline lineage and resolves the root cause (a feature window drift in PR~\#847) within hours. \textbf{Case~2} (right): an A/B experiment with hidden feature leakage---vibe coding approves a statistically significant but causally invalid lift, while $C$'s causal DAG enforces a feature parity contract that blocks rollout before inference begins.}
  \label{fig:democases}
  \vspace{-16pt}
\end{figure*}

Evidence is not a passive record; it is a first-class attribute of the graph, linked directly to claims in $C$.
For example, an edge ``dataset $D$ conforms to schema $S$'' should have pointers to validation jobs and their timestamps.
This auditability requires a precise correspondence between graph elements and artifacts:
\begin{itemize}[leftmargin=*, itemsep=2pt]
  \item \textit{Structural mapping:} nodes map to code units (packages, classes, functions), configuration blocks, or dataset definitions.
  \item \textit{Behavioral mapping:} edges map to explicit call sites, API contracts, data lineage, or testable invariants.
  \item \textit{Round-trip integrity:} changes in $C$ compile down to edits in $A$; changes in $A$ rehydrate back into $C$ via static analysis, runtime traces, and provenance logs.
\end{itemize}

\subsubsection{$C$ as an Operable Sandbox}
$C$ is where agents reason about structure: manipulating dependencies, running analyses, and evaluating counterfactuals without modifying artifacts.

For instance, the system should answer questions of the form:
\begin{quote}
If I delete this node in the consensus graph, which dependencies will break in the downstream code?
\end{quote}

Operationally, this requires that $C$ encodes enough dependency information (and links to evidence) to support ``what-if'' queries and rollback-safe planning.
If the answer depends on unknown or implicit dependencies, that ambiguity should appear as increased consensus entropy rather than as a silent guess.
$C$ also supports context-aware view synthesis: task-relevant subgraphs are projected on demand rather than exposing full topology, keeping $C$ cognitively tractable at scale.

\subsection{Round-Trip Synchronization Requirements}
Structural commitments in $C$ are useful only if they stay in correspondence with artifacts. This requires synchronization to be a first-class system property, not an optional documentation step.

\subsubsection{Synchronization as a First-Class Property}
A session alternates between consensus moves (proposing or modifying $C$) and realization moves (editing $A$ and rehydrating back into $C$). Synchronization operators couple these moves, preserving round-trip integrity and updating evidence links.
We posit two principal operators:
\begin{itemize}[leftmargin=*, itemsep=2pt]
  \item $\Phi: (C, A) \rightarrow A'$ \quad (\emph{realize}): given the current artifact baseline $A$, compile a proposed consensus change into incremental diffs.
  \item $\Psi: (A, C) \rightarrow C'$ \quad (\emph{rehydrate}): infer/update consensus structure from observed artifact diffs.
\end{itemize}
In practice, $\Psi$ is implemented via static analysis (e.g., AST) \cite{yamaguchi2014cpg}, dynamic invariant detection \cite{ernst2007daikon}, data lineage/provenance, and test results. These operators must run continuously, not as on-demand documentation scripts.

\subsubsection{Round-Trip Consistency as a Monitored Approximation}
Ideally, $\Psi(\Phi(C, A)) \approx C$: realizing a consensus change and rehydrating the result should recover the original structure.
In practice, strict equality is unattainable because $\Phi$ may discard structural information that has no artifact-level counterpart, and $\Psi$ is underdetermined.
We therefore treat round-trip consistency as a \emph{soft invariant}: the system continuously measures structural drift $d(C, \Psi(\Phi(C,A)))$ and flags divergence above a threshold, rather than demanding exact isomorphism.
This relaxation is essential for incremental adoption: partial consensus layers provide value before full round-trip fidelity.

When ambiguous, $\Psi$ produces ranked candidate deltas; the system runs cheap discriminating checks (type/schema checks, unit tests, static analyses) and escalates to human input when uncertainty remains high, increasing $\mathcal{H}(C\mid I)$ rather than making a silent guess. Tuning this threshold remains open, as it must balance the cost of interrupting the developer against the risk of silently misattributing intent. 

\subsubsection{Human-Gated Review as a Control Surface}
Every proposed $\Delta C$ arrives with predicted $\Delta A$ and attached evidence; the human decision is gated on structure, not code.
The critical operation is commit/rollback: $\Delta A$ is available for audit and trust calibration but is not the primary review surface.

\subsection{Autonomous Agentic Orchestration}
No single model can maintain $C$ reliably across the full scope of a project. Specialization is necessary: different agents must own different responsibilities over $C$, and their coordination must itself be mediated through $C$ rather than through opaque message exchange.

\subsubsection{Coordinated Multi-Agent Pipelines}
We propose a coordinated pipeline of specialized agents that operate over $C$ and $A$.
This decomposition is informed by emerging multi-agent software-development systems, while shifting the locus of control to a governable consensus artifact $C$ \cite{hong2024metagpt,qian2024chatdev,li2025aidev,wang2025genmentor}.
A representative decomposition uses four roles: an Architect translates intent into structure; a Builder implements artifacts; an Auditor generates evidence and detects conflicts; a Navigator projects the slice of $C$ relevant to the current task. Each agent operates over $C$ as the shared coordination surface rather than through message exchange alone.
Commit/Rollback remains a human decision in all cases.

\subsubsection{Agent--Agent Consensus}
Agents negotiate using richer internal representations, but every outcome is projected back into $C$ before commitment. Conflicts surface as competing $\Delta C$ alternatives with natural-language rationale, highlighted affected paths, and predicted test outcomes; humans commit or reject.

\subsubsection{Expert-Guided Evolution}
Agents that mine interaction traces to discover structural patterns can progressively reduce human intervention distance. 
If recurring corrections cluster around predictable missing relations, the system learns to make those relations first-class in $C$.
This transforms $C$ from a static record into a living knowledge base that grows more accurate as the project evolves and human feedback accumulates.
